\chapter{Range cores and low-rank interfaces}\label{ch:range}
\chapterlead{After storage and propagation have been organized, the remaining question is
which subspace has actually been resolved. Local ranks govern the width of
phase messages. Orthogonal factor updates carry that information into a small
core with the same nonzero spectrum when the full range has been captured.
An incomplete range requires an additional certificate, developed next.}
\section{Phase-local rank}\label{sec:final-18}

Let \(R\) be a newly generated Krylov residual and
\[
R_h=P_hR.
\]

\begin{theorem}[Minimal exact phase width]\label{claim:final-18-1}
Let
\[
r_h=\operatorname{rank}(R_h).
\]
Every phase-supported basis spanning \(\operatorname{range}(R_h)\) has at
least \(r_h\) vectors, and a rank-revealing QR or compact SVD attains exactly
\(r_h\).
\end{theorem}

This justifies nonuniform Krylov block widths by phase.


\begin{figure}[htbp]
\centering
\includegraphics[width=.82\linewidth]{figure_05_phase_local_rank.png}
\caption{Phase-local rank as a dimension law. The archived six-phase audit records constructed rank, numerical rank and minimum exact width of one through six, with the remaining capacity wasted by a uniform width of twelve. A rank-revealing factorization keeps the new directions. This summary replaces an annotated singular-value curve whose raw plotted vector was not archived.}
\label{fig:phase-local-rank}
\end{figure}


\section{Low-rank interface messages}\label{sec:final-19}

If an inter-phase coupling has exact factorization
\[
P_hHP_g
=
U_{hg}V_{hg}^\top
\]
of rank \(r_{hg}\), then
\[
\boxed{
P_hHP_gx
=
U_{hg}(V_{hg}^\top x).
}
\]

This representation reduces storage and communication.  It does not
automatically reduce wall-clock time unless the factors are applied through
batched or fused kernels. A Python reference implementation can lose to
dense BLAS despite lower arithmetic and storage.


\section{Exact QR-core updates}\label{sec:final-20}

Suppose
\[
A=QR,
\qquad
Q^\top Q=I,
\]
and append a new block \(U\).  Set
\[
C=Q^\top U,
\qquad
F=U-QC,
\qquad
F=Q_\perp R_\perp,
\qquad Q_\perp^\top Q_\perp=I,
\qquad Q^\top Q_\perp=0.
\]
Here the columns of \(Q_\perp\) form an orthonormal basis of the nonzero
residual range. Zero residual directions are omitted.

\begin{theorem}[Exact QR-core update]\label{claim:final-20-1}
\[
[A,U]
=
[Q,Q_\perp]
\begin{pmatrix}
R&C\\
0&R_\perp
\end{pmatrix}.
\]
Therefore the nonzero eigenvalues of
\([A,U][A,U]^\top\) are exactly the nonzero eigenvalues of the small
core Gram, with multiplicity. If $\operatorname{range}(Q)=\operatorname{range}(A)$,
the compact residual update also spans exactly $\operatorname{range}([A,U])$.
In that case the core has full row rank and its Gram matrix has no additional
zero eigenvalues.
\end{theorem}

\begin{proof}
The first block product is \(QR=A\). The second is
\(QC+Q_\perp R_\perp=QC+F=U\).  The stated orthogonality
conditions make \([Q,Q_\perp]\) an isometry on the compact core
space.  Left multiplication by this isometry preserves all nonzero singular
values of the core, and squaring them yields the nonzero Gram eigenvalues.
If $Q$ spans exactly the old range and $Q_\perp$ spans exactly the residual
range, their combined range is precisely that of the appended factor.
The core then has full row rank. Without this exact-range hypothesis, unused
basis directions can introduce zero eigenvalues into the core Gram.
\end{proof}

This is useful for overlapping genomic windows and small incremental variant
updates.


\begin{figure}[htbp]
\centering
\includegraphics[width=.84\linewidth]{figure_06_qr_core_equivalence.png}
\caption{QR-core spectral equivalence under the stated orthonormality hypotheses. The diagram pairs the exact factor identity with the archived maximum absolute and relative \(\ell_2\) singular-spectrum discrepancies for a 128-row factor with 180 old and 37 appended columns. Recorded audit metrics replace the overlaid singular-value curves, whose raw vectors were not archived.}
\label{fig:qr-core-equivalence}
\end{figure}


\section{Range-core spectral compression}\label{sec:final-21}

Suppose an orthonormal basis \(Q_r\) captures the factor range:
\[
\operatorname{range}(A)\subseteq\operatorname{range}(Q_r).
\]
Then
\[
A=Q_rB,
\qquad
B=Q_r^\top A.
\]

\begin{theorem}[Exact range core]\label{claim:final-21-1}
The nonzero eigenvalues of
\[
K=AA^\top
\]
are exactly the nonzero eigenvalues of
\[
BB^\top,
\]
with multiplicity. If $\operatorname{range}(Q_r)=\operatorname{range}(A)$,
then $BB^\top$ is positive definite on the core space and its entire spectrum
is the nonzero spectrum of $K$.
\end{theorem}

\begin{proof}
Since $A=Q_rB$ and $Q_r^\top Q_r=I$, one has
$AA^\top=Q_r(BB^\top)Q_r^\top$. The map $x\mapsto Q_rx$ is an isometry
from the core space onto $\operatorname{range}(Q_r)$. On that range the two
operators are conjugate, and $AA^\top$ vanishes on its orthogonal complement.
Hence their nonzero eigenvalues agree with multiplicity. If the captured
range equals $\operatorname{range}(A)$, then $B$ has full row rank, so
$BB^\top$ is positive definite. Without that equality, the core may retain
additional zero eigenvalues.
\end{proof}


Thus the computational pipeline can terminate its data-streaming stage after
the range is captured:
\[
\boxed{
\text{streamed data}
\to
\text{small spectral core}
\to
\text{tail evaluation}.
}
\]



\clearpage
\section{An exact core calculation}
Consider the orthonormal range basis and core factor
\[
U=\begin{pmatrix}1&0\\0&1/\sqrt2\\0&1/\sqrt2\end{pmatrix},
\qquad
B=\begin{pmatrix}1&0&1\\0&2&0\end{pmatrix},
\qquad A=UB.
\]
Direct multiplication gives
\[
BB^\top=\begin{pmatrix}2&0\\0&4\end{pmatrix},
\qquad
AA^\top=\begin{pmatrix}2&0&0\\0&2&2\\0&2&2\end{pmatrix}.
\]
The large operator has eigenvalues $4,2,0$. Its two-dimensional core has
exactly the nonzero eigenvalues $4,2$. The Gaussian quadratic form therefore
has the law $4Z_1^2+2Z_2^2$, whether it is computed from the full factor, the
full Gram matrix, or the core.

\subsection*{Why the nonzero-spectrum qualification matters}
Take instead
\[
A=\begin{pmatrix}1&0\\0&0\end{pmatrix},\qquad Q_r=I_2,\qquad B=Q_r^\top A=A.
\]
The captured subspace contains the range of $A$, but is larger than that range.
Both Gram matrices have spectrum $\{1,0\}$. Their nonzero spectrum is
$\{1\}$. Thus the nonzero spectrum of $AA^\top$ cannot be identified with the
\emph{entire} spectrum of $BB^\top$. Choosing the compact basis
$Q_r=(1,0)^\top$ instead gives $B=(1,0)$ and the one-dimensional core Gram
$(1)$. The extra zero records an unused basis direction, not an additional
random contribution. This is why range containment and exact-range equality
have different conclusions in Theorem~\ref{claim:final-21-1}.

The first example in this section has active rank two. It separates algebraic
exactness from backend coverage. The spectral reduction and analytic packet theory remain
valid, but the implemented directed positive-tail backend described in
Chapter~\ref{ch:intervals} requires active rank greater than two. Capturing a
range exactly does not remove a later numerical method's hypotheses.
